\documentclass[11pt]{article}

\usepackage[margin=1in]{geometry}
\usepackage{hyperref}

\title{Nakafa Lab: Toward Affordable Disease-Modifying Research for Thalassemia}
\author{Nakafa Lab}
\date{Draft started 2026-04-26}

\begin{document}
\maketitle

\begin{abstract}
This draft defines the first Nakafa Lab research program: a source-backed,
AI-assisted search for credible and affordable thalassemia treatments. The work
combines clinical literature review, mechanism mapping, computational
hypothesis generation, Indonesia access analysis, and a separate Islamic
research lane. This document is not medical advice.
\end{abstract}

\section{Problem}

Thalassemia remains a major inherited blood disorder. For severe
beta-thalassemia, standard management can include lifelong transfusion and iron
chelation. Curative or disease-modifying therapies exist, but cost,
infrastructure, and eligibility limit access \cite{farmakis2022tif}.

\section{Research Objective}

The objective is to identify credible treatment paths that can reduce disease
burden and become accessible to patients who cannot realistically access
high-cost curative therapies.

\section{Initial Evidence Map}

The initial evidence map covers clinical guidelines, luspatercept, mitapivat,
gene addition therapy, CRISPR-edited autologous stem-cell therapy, fetal
hemoglobin induction, erythropoiesis, iron overload, and Indonesia-specific
access constraints \cite{cappellini2020luspatercept,taher2025energize,locatelli2024exaCel,fda2025aqvesme}.

\section{Islamic Research Lane}

The Islamic research lane begins with Quran 16:68-69 and expands through
translation, tafsir, sirah, hadith, and scholarly review. This lane can generate
ethical framing and hypotheses, but biological claims must pass scientific
validation \cite{quran1668}.

\section{Safety}

No output from this project should be used for diagnosis, dosing, or treatment
without licensed clinical review.

\bibliographystyle{plain}
\bibliography{references}

\end{document}
